\section{Preliminaries}
\subsection{Multi-Agent Debate}
Multi-agent debate (MAD) is an emergent paradigm for collective reasoning, where multiple language model agents iteratively exchange their intermediate thoughts and conclusions in order to reach a more accurate and robust consensus than any single agent alone. A central research focus within MAD is the design of communication topologies—that is, how agents share information and interact across debate rounds.

Early studies typically employed fully connected topologies, allowing all agents to access each other's responses at every round. However, such dense communication can be inefficient and costly, particularly as the number of agents grows. To address this, subsequent works have explored various forms of sparse and structured topologies. For instance, S-MAD \citep{li2024improving} adopts fixed sparse patterns such as rings (only adjacent agents are visible) or stars (only centric agent are visible), enabling each agent to communicate only with a subset of neighbors, thereby reducing token consumption. Other approaches, like Group Debate (GD) \citep{liu2024groupdebate}, partition agents into static groups, which debate internally before group-level aggregation. Similarly, two-stage topologies such as S$^2$-MAD \citep{zeng2025s} first organize local debates within small groups and then aggregate results through a secondary global round if group opinions differ. These strategies all aim to balance the trade-off consensus formation, and computational efficiency.

\subsection{Static Workflow Generation and Dynamic Topology Control}

\newtext{Multi-agent debate (MAD) systems rely on efficient inter-agent communication topologies. This challenge shares many similarities with recent research about automatic generation of agent workflows. such as MACNET \citep{qian2025ScalingLargeLanguage}, AFlow \citep{zhang2025AFlowAutomatingAgentic}, and MaAS \citep{zhang2025MultiagentArchitectureSearch}, which primarily focus on optimizing the static workflow or agentic pipeline. For example, MACNET analyzes the performance of predefined static topologies (e.g., chain, mesh) . AFlow uses MCTS to search for a single, optimal static workflow for a given task domain. MaAS refines this by using a controller to select a single, query-dependent static workflow at the beginning of a task. Despite above advances, however, most existing multi-agent LLM frameworks rely on static or manually specified topologies, which cannot adapt to the dynamic needs or varying complexity of different tasks.}

This motivates the development of adaptive and principled methods for controlling agent interactions. \newtext{Reinforcement learning (RL) is an end-to-end solution which is capable of both modeling and scheduling temporal dynamics during debate. Hierarchical RL frameworks like FMH \citep{ahilan2019FeudalMultiAgentHierarchies} use multi-faceted rewards, where the objective is semantic task decomposition via manager-worker subgoaling.}

\newtext{Our work addresses a fundamentally different problem. RUMAD is not concerned with generating a static, task-level workflow. Instead, we focus on dynamic, round-by-round communication management within a debate. Our RL controller is invoked at each step to adaptively reconfigure the communication graph based on the evolving state of the debate (e.g., agent agreement, semantic similarity). This allows RUMAD to achieve fine-grained efficiency, such as pruning connections to agents that have already reached consensus, a capability not present in static workflow models. Considering hierarchical reward modeling, RUMAD's multi-objective reward function is novel in its content-agnostic nature, balancing the system-level trade-offs between solution quality, group cohesion, and computational cost.}


\subsection{Markov Decision Process}
A Markov Decision Process (MDP) is a mathematical framework for modeling sequential decision-making problems, where an agent interacts with an environment over discrete time steps. 
% Formally, an MDP is defined by a tuple $(\mathcal{S}, \mathcal{A}, P, r, \gamma)$, where $\mathcal{S}$ denotes the state space, $\mathcal{A}$ the action space, $P$ the state transition probability, $r$ the reward function, and $\gamma$ the discount factor. 
At each step, the agent observes the current state, selects an action, and receives a reward and a new state from the environment. The objective is to find a policy that maximizes the expected cumulative reward \citep{suttonReinforcementLearningIntroduction2018}.
In our multi-agent debate (MAD) setting, this framework allows the controller to base its topology decisions solely on the observable debate state at each round, ensuring a neutral and adaptive decision process that does not require access to agents’ internal semantic content.


\subsection{Proximal Policy Optimization}
Proximal Policy Optimization (PPO) \citep{schulmanProximalPolicyOptimization2017a} is a widely used reinforcement learning algorithm for training policies in high-dimensional and dynamic environments. PPO employs an actor-critic framework, where a policy network proposes actions and a value network estimates the expected return for each state. By optimizing a clipped surrogate objective, PPO achieves stable and efficient policy updates, and has demonstrated robust performance across various sequential decision-making tasks.
In our MAD framework, the value network provides reliable estimates of the debate state’s long-term potential, which in turn stabilizes and guides the learning of RUMAD about adaptive topology control strategies. Additionally, we design a multi-objective reward model to capture the diverse goals of the debate process, and employ Generalized Advantage Estimation (GAE) \citep{schulmanHighDimensionalContinuousControl2018} to further improve learning efficiency and convergence.

